\MulaiModul{02}{Masukan}

\subsection{Tujuan}

\begin{tujuan}
Setelah menyelesaikan modul ini, mahasiswa mampu:
\begin{enumerate}
    \item membaca satu karakter dengan \texttt{INT 21h} fungsi
          \texttt{AH=01h} dan menjelaskan hasilnya di \texttt{AL};
    \item membaca satu baris dengan \texttt{INT 21h} fungsi
          \texttt{AH=0Ah} ke dalam penyangga;
    \item menjelaskan susunan byte penyangga fungsi \texttt{0Ah};
    \item mencetak kembali isi masukan dengan layanan \texttt{AH=02h};
    \item merakit, menautkan, dan menjalankan kedua program dengan
          TASM dan TLINK.
\end{enumerate}
\end{tujuan}

\subsection{Dasar singkat}

\subsubsection*{Layanan 01h: membaca satu karakter}

Layanan \texttt{AH=01h} menunggu satu tombol ditekan, lalu menaruh
kode ASCII tombol itu pada \texttt{AL}. Karakter tersebut juga
ditampilkan ke layar; sifat ini disebut \emph{echo}. Karena hasilnya
berada di \texttt{AL}, program dapat memindahkannya ke register lain
atau mencetaknya kembali. Satu panggilan layanan ini hanya membaca
satu karakter.

\subsubsection*{Layanan 0Ah: membaca satu baris}

Layanan \texttt{AH=0Ah} membaca beberapa karakter sekaligus sampai
tombol Enter ditekan. Sebelum \texttt{INT 21h}, \texttt{DS:DX} harus
menunjuk awal penyangga. Penyangga memiliki susunan tiga bagian
seperti ditunjukkan pada Tabel~\ref{tab:m02-buffer}.

\begin{table}[H]
    \centering
    \caption{Susunan penyangga layanan \texttt{AH=0Ah}.}
    \label{tab:m02-buffer}
    \small
    \begin{tabular}{@{}lll@{}}
        \toprule
        \textbf{Byte} & \textbf{Diisi oleh} & \textbf{Isi} \\
        \midrule
        0 & program & kapasitas, termasuk tombol Enter \\
        1 & DOS & jumlah karakter yang diketik, tanpa Enter \\
        2 dan seterusnya & DOS & karakter masukan, diakhiri \texttt{0Dh} \\
        \bottomrule
    \end{tabular}
\end{table}

Karena satu byte disediakan untuk tombol Enter, jumlah karakter yang
dapat diketik paling banyak adalah kapasitas dikurangi satu. Bila
kapasitas bernilai 40, pengguna dapat mengetik sampai 39 karakter.
Angka pada byte 1 diisi DOS setelah Enter ditekan dan tidak
memasukkan \texttt{0Dh}. Isi masukan dimulai pada byte 2 dan diakhiri
\texttt{0Dh}. Dengan demikian, byte 1 dapat dipakai sebagai pencacah
untuk mencetak tepat sebanyak karakter yang diketik.

\subsubsection*{Mencetak kembali isi penyangga}

Isi penyangga dicetak kembali dengan layanan \texttt{AH=02h}.
Karakter diambil satu per satu dari byte 2, lalu \texttt{DL} diisi
karakter itu dan \texttt{INT 21h} dijalankan. Pencacah diambil dari
byte 1, sehingga \texttt{0Dh} pada akhir masukan tidak ikut tercetak.

Pengulangan kecil yang dipakai di sini hanya alat bantu untuk
menampilkan masukan. Bagian ini bukan materi perulangan dan tidak
memakai instruksi \texttt{LOOP}. Pembahasan perulangan sebagai topik
tersendiri diberikan pada modul berikutnya.

Tabel~\ref{tab:m02-layanan} merangkum layanan masukan dan keluaran
yang dipakai pada modul ini.

\begin{table}[H]
    \centering
    \caption{Layanan yang dipakai pada Modul 02.}
    \label{tab:m02-layanan}
    \small
    \begin{tabular}{@{}lll@{}}
        \toprule
        \textbf{Layanan} & \textbf{Register masukan} & \textbf{Hasil} \\
        \midrule
        \texttt{AH=01h} & -- & \texttt{AL} = kode ASCII, ikut tercetak \\
        \texttt{AH=0Ah} & \texttt{DS:DX} = awal penyangga &
        isi penyangga terisi \\
        \texttt{AH=02h} & \texttt{DL} = satu karakter &
        karakter tercetak \\
        \texttt{AH=4Ch} & \texttt{AL} = kode keluar &
        program berakhir \\
        \bottomrule
    \end{tabular}
\end{table}

\begin{rujukan}
Rujukan singkat: Gunawan, Bagian 6.8; S'to, Bab XV bagian masukan
\emph{keyboard}; buku ajar, \texttt{input-nama.asm} dan
\texttt{chapters/09}.
\end{rujukan}

\subsection{Alat dan bahan}

\begin{itemize}
    \item Komputer dengan emu8086 terpasang.
    \item DOSBox dengan TASM, TLINK, dan DEBUG.
    \item Berkas \texttt{code/m02-baca-karakter.asm} dan
          \texttt{code/m02-baca-baris.asm}.
    \item Lembar jurnal praktikum.
\end{itemize}

\subsection{Contoh program}

Program pertama membaca satu karakter dengan layanan \texttt{01h}.
Karakter yang diketik tampil sekali karena \emph{echo}, lalu dicetak
sekali lagi oleh layanan \texttt{02h} dari isi \texttt{AL}.

\lstinputlisting[
    caption={Membaca satu karakter dengan \texttt{AH=01h}.},
    label={lst:m02-karakter}
]{code/m02-baca-karakter.asm}

Program kedua membaca satu baris dengan layanan \texttt{0Ah}.
Setelah Enter ditekan, byte 1 berisi jumlah karakter. Program
mengambil jumlah itu ke \texttt{CX}, lalu mencetak isi mulai byte 2
sebanyak jumlah tersebut dengan layanan \texttt{02h}.

\lstinputlisting[
    caption={Membaca satu baris dengan \texttt{AH=0Ah}.},
    label={lst:m02-baris}
]{code/m02-baca-baris.asm}

\subsection{Langkah kerja}

\begin{enumerate}
    \item Siapkan folder kerja, lalu salin kedua berkas \texttt{.asm}
          ke dalamnya.
    \item Rakit dan tautkan program pertama. Jalankan kedua perintah
          berikut pada baris terpisah:
          \begin{perintah}
tasm m02-baca-karakter.asm\\
tlink m02-baca-karakter.obj
\end{perintah}
          Jalankan hasilnya, ketik satu karakter, dan catat keluarannya.
    \item Rakit dan tautkan program kedua dengan perintah yang sama,
          lalu jalankan. Ketik nama Anda (minimal satu karakter) dan
          tekan Enter. Catat keluarannya.
    \item Ulangi program kedua dengan kalimat yang lebih panjang dari
          kapasitas penyangga. Amati jumlah karakter yang tercetak
          dan jelaskan mengapa sisanya tidak terbaca.
    \item Periksa penyangga dengan DEBUG: jalankan
          \texttt{debug m02-baca-baris.exe}, ketik \texttt{r}, lalu
          \texttt{d ds:0 L 60}. Tandai byte kapasitas, byte jumlah,
          karakter masukan, dan byte \texttt{0Dh}.
    \item Ubah kapasitas penyangga dari 40 menjadi 10, rakit ulang,
          lalu ulangi langkah 3. Catat perbedaan jumlah karakter yang
          dapat diketik.
\end{enumerate}

\subsection{Tugas}

\begin{enumerate}
    \item Jelaskan perbedaan layanan \texttt{AH=01h} dan
          \texttt{AH=0Ah} dari segi jumlah karakter dan hasilnya.
    \item Jelaskan isi byte 0, byte 1, dan byte 2 pada penyangga
          fungsi \texttt{0Ah}.
    \item Jelaskan mengapa karakter pada program pertama tampil dua
          kali.
    \item Bila byte kapasitas bernilai 10, berapa karakter paling
          banyak yang dapat diketik sebelum Enter? Jelaskan alasannya.
    \item Tambahkan kalimat pembuka sebelum masukan pada program
          kedua dengan layanan \texttt{AH=09h}, lalu cetak ulang
          masukan seperti semula.
\end{enumerate}

\subsection{Hasil yang dikumpulkan}

\begin{enumerate}
    \item Kedua berkas \texttt{.asm} beserta versi modifikasinya.
    \item Tangkapan layar keluaran kedua program.
    \item Tabel isi penyangga hasil pemeriksaan DEBUG, termasuk byte
          kapasitas, byte jumlah, dan byte \texttt{0Dh}.
    \item Lembar jurnal berisi jawaban langkah kerja dan tugas.
\end{enumerate}
